\documentclass[11pt]{article}
\usepackage{amsmath,amssymb,geometry}
\geometry{margin=1in}
\newcommand{\Tr}{\operatorname{Tr}}
\begin{document}
\section*{Current result: the maximum is two}
The auxiliary-spin inertia theorem now proves $k_{\max}(13,2)=2$ for arbitrary
complex subspaces. With $X=\mathbf J\cdot\mathbf S$ and auxiliary spin $3/2$,
$W=X^2-9X/2-975I/16$ has eleven positive and forty-five negative eigenvalues,
and no zero eigenvalues. Its compression to the $4k$-dimensional code--auxiliary
space vanishes. Projection into the positive eigenspace is injective, proving
$4k\le11$. The exact two-dimensional family below attains this bound.
The complete standalone proof is in \texttt{inertia/proof.tex}.
\section*{Historical investigations, superseded by the global proof}
The former unresolved status and weaker conditional results below describe work
before the auxiliary-spin theorem. They are retained for reproducibility and
are not the current catalogue status.

\title{Thirteen symmetric qubits, second-order anticoherent subspaces: current bounds}
\author{Research note}
\date{}
\maketitle

\paragraph{Current rigorous conclusion.}
The maximal dimension satisfies
\[
 \boxed{2\le k_{\max}(13,2)\le3.}
\]
The conjecture that the maximum is two has not been proved by the calculations recorded here. No three-dimensional counterexample has been constructed either.

\paragraph{An exact two-dimensional witness.}
An orthonormal basis is
\[
 |u\rangle=\frac{\sqrt{10}|D_{13}^{(0)}\rangle+\sqrt{26}|D_{13}^{(9)}\rangle}{6},
 \qquad
 |v\rangle=\frac{-\sqrt{26}|D_{13}^{(4)}\rangle+\sqrt{10}|D_{13}^{(13)}\rangle}{6}.
\]
All Dicke support gaps within and between these vectors exceed two. The first moments vanish and each second diagonal moment equals
$j(j+1)/3=65/4$, with $j=13/2$. More directly, exact Dicke splitting verifies every entry of
\[
 \Tr_{11}(|u_a\rangle\langle u_b|)=\delta_{ab}\frac{I_3}{3}.
\]
All four ordered basis pairs, comprising 36 matrix entries, pass this identity. This certifies the entire span, rather than only its basis states.

\paragraph{Why the upper bound is three.}
For a second-order AC subspace of dimension $k$, the vectors obtained from an orthonormal basis by the four operators
\[
 I,\quad J_x/\sqrt{65/4},\quad J_y/\sqrt{65/4},\quad J_z/\sqrt{65/4}
\]
are mutually orthonormal. The symmetric Hilbert space has dimension fourteen, so $4k\le14$ and $k\le3$.

\paragraph{What a hypothetical dimension-three subspace would imply.}
Let $P$ be its orthogonal projector and write
$\Phi(A)=\sum_aJ_aAJ_a$. The unused spin-error space is represented by
\[
 Q=I-P-\frac4{65}\Phi(P).
\]
It is an orthogonal projector of rank two, and
\[
 PQ=0,\qquad PJ_aQ=0.
\]
For the irreducible tensor components of the operators,
\[
 P_L=\beta_LQ_L,\qquad
 \beta_0=\frac32,\qquad
 \beta_L=\frac{65}{2L(L+1)-260}\quad(1\le L\le13).
\]
The first two multipoles of $Q$ vanish. With
$a_L=\|Q_L\|_{\rm HS}^2$, useful necessary scalar identities are
\[
 a_0=\frac27,\quad a_1=a_2=0,\quad
 \sum_La_L=2,\quad
 \sum_L\beta_La_L=0,\quad
 \sum_L\beta_L^2a_L=3.
\]
The last two equalities are $\Tr(PQ)=0$ and $\Tr(P^2)=3$.

The complementary Schmidt space gives another projector,
\[
 H=I_{\operatorname{Sym}^{11}}-3\Tr_2P,
 \qquad \operatorname{rank}H=3.
\]
Its existence follows from the nine mutually orthogonal Schmidt partner vectors of the three code states across the two-qubit versus eleven-qubit split.

\paragraph{Averaged anticoherence is not a code condition.}
Vanishing $Q_1,Q_2$ means that the mixed state $Q/2$ has the required averaged first and second moments. It does not imply $QT_{Lq}Q=0$. Thus one cannot treat the rank-two complement as a second-order AC subspace. For a simple illustration in spin three, the equally weighted mixture of $|m=2\rangle$ and $|m=-2\rangle$ is second-order AC, but the two-dimensional span is not even first-order AC: its compressed $J_z$ has eigenvalues $2$ and $-2$.

\paragraph{What the relaxations establish, and what they do not.}
An exact rational point satisfies the tested scalar-invariant relaxation, including 450 nonnegativity inequalities and nine equality rows. These incorporate purity, spin-tensor transitions, two-copy probabilities, the spin-error and Schmidt complements, and mixed-spin positivity with time reversal. The point, expressed as the squared multipole lengths of the hypothetical rank-three projector $P$, is saved in \texttt{relaxation-report.json}; \texttt{verify\_scalar.py} reconstructs and verifies its constraints exactly. Its existence proves that these scalar constraints alone cannot exclude dimension three.

Stronger matrix-moment relaxations were also investigated. They retain some correlations discarded by scalar multipole lengths, using exterior products and rotation-averaged copy moments. Feasibility of such a relaxation does not produce a physical subspace: a positive moment matrix need not come from an actual rank-three AC projector, or factor into the required identical-copy pure moments. Conversely, a numerical infeasibility or matrix-rank report is insufficient without an exact certificate and an independently checked representation.

Nine direct searches over complex bases, including seeds derived from the known fourteen-qubit construction, found no exact zero of the whole-subspace constraints. This is numerical evidence only. The strongest tested matrix relaxation, including partial-transpose positivity and the full wedge-copy partial transpose, remained numerically feasible after its known zero directions were removed. Its smallest remaining block eigenvalue was approximately $9.50\times10^{-6}$; no physical three-dimensional subspace was recovered from it.

Before the auxiliary-spin proof, the outcome was the rigorous interval $2\le k_{\max}(13,2)\le3$. Settling it requires either an exact exclusion covering all complex three-dimensional subspaces or an explicit, fully verified dimension-three construction.
\section*{A continuous family and a rigorous nonextension result}
Write $D_r=|D_{13}^{(r)}\rangle$. For every real $\phi$, set
\[
 u=\frac{\sqrt{10}D_0+\sqrt{26}D_9}{6},\qquad
 v_\phi=\frac{\sqrt{26}D_4+e^{i\phi}\sqrt{10}D_{13}}6.
\]
The two vectors have disjoint supports and unit norms. Their support union
$\{0,4,9,13\}$ has all gaps at least four. Consequently all off-diagonal
spin-tensor matrix elements through rank two vanish. Each vector has
$\langle J_z\rangle=0$ and $\langle J_z^2\rangle=65/4$; cross diagonal
elements vanish by disjoint support. These identities prove that the entire
span is second-order AC, for every $\phi$. It is not third-order AC:
$\langle u|J_z^3|u\rangle=65\ne0$.

\paragraph{Inequivalence under collective rotations.}
Let $P_\phi$ be the projector onto this span, and define spin time reversal
by $\Theta D_r=(-1)^rD_{13-r}$, with antilinear extension. Direct calculation gives
\[
 \operatorname{Tr}(P_\phi\Theta P_\phi\Theta^{-1})
 =2-\frac{65}{81}(1+\cos\phi).
\]
Time reversal commutes with collective rotations, so this quantity is invariant.
It increases strictly from $32/81$ to $2$ as $\phi$ runs from $0$ to $\pi$.
Thus these parameters give pairwise rotation-inequivalent subspaces.

\paragraph{Every member is inclusion-maximal.}
Put $z=e^{i\phi}$ and use $J_zD_r=(13/2-r)D_r$ and
$J_-D_r=\sqrt{(r+1)(13-r)}D_{r+1}$. Form the matrix with columns
\begin{align*}
 &u,J_zu,v_\phi,J_zv_\phi,J_+u,J_+^2u,J_-u,\{J_z,J_-\}u,\\
 &J_-v_\phi,J_-^2v_\phi,J_+v_\phi,\{J_z,J_+\}v_\phi,J_-^2u,J_+^2v_\phi.
\end{align*}
In the normalized Dicke basis its determinant is exactly
\[
 -2413404500000\,z^2(z-11),
\]
which is nonzero when $|z|=1$. Every column is an image under an operator
of tensor rank at most two. A vector extending the span to a larger second-order
AC subspace would be orthogonal to all fourteen columns, hence zero.
This proves inclusion-maximality, not a global maximum-dimension theorem.
The symbolic verifier checks all 36 reduced-state matrix entries, the determinant,
the time-reversal invariant, and the nonzero third moment exactly.
\section*{A stronger conditional exclusion}
An exact positive certificate also excludes every three-dimensional second-order
AC code containing $u$ itself, regardless of the other two basis vectors.
Their orthogonality to all rank-at-most-two tensor images of $u$ places them in
\[
\mathcal K=\operatorname{span}\left\{(\sqrt{11}D_2-D_{11})/\sqrt{12},
D_3,D_4,D_5,D_6,D_{12},D_{13}\right\}.
\]
The ninefold axial rotation preserves $u$ projectively and preserves $\mathcal K$.
Averaging the exterior-product density of the other two vectors gives a positive
operator on $\Lambda^2\mathcal K$, with 48 real invariant block coordinates.
Forty independent exact equations leave an eight-dimensional kernel. A strictly
positive integer identity on all 21 block diagonals annihilates this kernel,
forcing the operator to vanish and contradicting its nonzero trace.
The complete proof, certificate and two physical reconstructions are supplied
in the accompanying \texttt{containing-u/} directory.

Every second-order AC state with at most two Dicke components, along any axis,
is a collective rotation of $u$. Indeed its opposite magnetic labels $a,-b$
satisfy $ab=65/4$. Both $2a$ and $2b$ are positive odd integers at most 13,
so their unordered pair is $\{5,13\}$. The weights are fixed by zero mean,
and axial rotation removes the relative phase. Consequently any hypothetical
dimension-three code must avoid every such two-component state.
This remains conditional: no assertion that all three-dimensional codes contain
such a state is made, and this conditional argument alone left the interval $2$--$3$; the auxiliary-spin theorem now closes that gap.
\end{document}
